\chapter{与 C 互操作}
\label{ch:21}

到第 20 章为止，你的程序都活在自己的世界里。可现实中有太多东西是现成的 C 库：压缩、加密、数据库、图形、系统调用……你不可能全都重写一遍。

这一章讲怎么跨过 ChaLi 和 C 的边界，两个方向：

\begin{itemize}

\item ChaLi 调 C——用 \texttt{\#abi} 指令调动态库里导出的函数；

\item C 调 ChaLi——把你的函数交给 C 的运行时。

\end{itemize}

\section{三种集成方式}

工具链支持三条路，按「需要多少 C 侧配合」排：

\begin{center}
\begin{tabular}{lll}
\toprule
方式 & 用在哪 & 你要准备什么 \\
\midrule
\texttt{\#abi} 指令 & 调动态库里按 C ABI 导出的函数 & 一份 \texttt{.\allowbreak{}dll} / \texttt{.\allowbreak{}so} \\
宿主胶水绑定 & 把 C++ 函数注册进工具链，ChaLi 侧当普通函数调 & 写一点 C++ \\
原生库嵌入 & 宿主调用 AOT 导出的 C ABI & 原生动态库和生成的头文件 \\
\bottomrule
\end{tabular}
\end{center}

\texttt{\#abi} 指令是最轻的一条——不用改工具链，不用重新编译，运行时直接加载库。

\section{\texttt{\#abi} 指令：调动态库里的 C 函数}

假设你有一个 C 文件 \texttt{native.\allowbreak{}c}：

\par
\medskip
\begin{lstlisting}[style=c]
#include <stdint.h>

#if defined(_WIN32)
#define CHALI_ABI_EXPORT __declspec(dllexport)
#else
#define CHALI_ABI_EXPORT
#endif

CHALI_ABI_EXPORT int32_t abi_add(int32_t left, int32_t right) {
    return left + right;
}

CHALI_ABI_EXPORT const char *abi_name(void) {
    return "native";
}
\end{lstlisting}
\par
\medskip

先把它编成动态库（说句实话：「怎么编动态库」是 C 的事，和 ChaLi 无关）：

\par
\medskip
\begin{lstlisting}[style=sh]
gcc -shared -o native.dll native.c
\end{lstlisting}
\par
\medskip

然后在 ChaLi 里声明想用的函数：

\par
\medskip
\begin{lstlisting}[style=chali]
#abi("C", "./native.dll", "abi_add")
int32 native_add(int32 left, int32 right);
#abi("C", "./native.dll", "abi_name")
String native_name();

void main()
{
    println("{} {}", native_add((int32)20, (int32)22), native_name());
}
\end{lstlisting}
\par
\medskip

\par
\medskip
\begin{lstlisting}[style=out]
42 native
\end{lstlisting}
\par
\medskip

导入使用逐函数编译器指令：

\par
\medskip
\begin{lstlisting}[style=out]
#abi("C", "<动态库路径>", "<库里的符号名>")
<返回类型> <ChaLi侧名字>(<形参表>);
\end{lstlisting}
\par
\medskip

第三个参数是库里真实导出的符号名，函数声明使用的是 ChaLi 中的本地名称。两者可以不一样——给一批 C 函数起统一风格的 ChaLi 名，就靠它。

\begin{tip}{提示}
路径是运行时按当前工作目录找的。所以 \texttt{.\allowbreak{}/\allowbreak{}native.\allowbreak{}dll} 得跟你的程序在同一个目录下跑。
\end{tip}

\section{结构体、引用、字符串}

标量只是热身。真实场景里你还要传结构体（按值）、引用（原地改）、字符串：

\par
\medskip
\begin{lstlisting}[style=c]
API Point move_point(Point p, int32_t dx, int32_t dy) {
    p.x += dx; p.y += dy;
    return p;
}
API void bump(int32_t *value) { ++*value; }
API const char *greet(void) { return "hello from C"; }
\end{lstlisting}
\par
\medskip

ChaLi 侧：

\par
\medskip
\begin{lstlisting}[style=chali]
struct Point
{
    int32 x;
    int32 y;
};

#abi("C", "./abi_demo.dll", "move_point")
Point move_point(Point p, int32 dx, int32 dy);
#abi("C", "./abi_demo.dll", "bump")
void bump(int32 mut& value);
#abi("C", "./abi_demo.dll", "greet")
String greet();

void main()
{
    Point p;
    p.x = 1;
    p.y = 2;
    Point q = move_point(p, (int32)10, (int32)20);
    println("{} {} / {} {}", p.x, p.y, q.x, q.y);

    int32 n = 41;
    bump(&n);
    println("{}", n);

    println("{}", greet());
}
\end{lstlisting}
\par
\medskip

\par
\medskip
\begin{lstlisting}[style=out]
1 2 / 11 22
42
hello from C
\end{lstlisting}
\par
\medskip

三个要点：

\begin{enumerate}

\item 结构体按值传——C 的 \texttt{Point p} 对应 ChaLi 的 \texttt{Point p}，布局默认按 C 规则对齐。所以 ChaLi 这边得声明一个字段顺序、类型都一致的 \texttt{struct}；

\item C 的 \texttt{int32\_t *value} 对应 \texttt{int32 mut\& value}（第 7 章的可变引用），调用处写 \texttt{bump(\allowbreak{}\&n)\allowbreak{}}——取地址用 \texttt{\&}，和 ChaLi 内部引用是同一套写法；

\item C 的 \texttt{const char *} 对应 \texttt{String}。

\end{enumerate}

\subsection*{支持的类型一览}

{\small\setlength{\tabcolsep}{4pt}
\begin{longtable}{>{\raggedright\arraybackslash}p{\dimexpr 0.586\linewidth-2\tabcolsep\relax}>{\raggedright\arraybackslash}p{\dimexpr 0.394\linewidth-2\tabcolsep\relax}}
\toprule
ChaLi 写法 & 含义 \\
\midrule
\endhead
\texttt{int8}…\texttt{int64}、\texttt{uint8}…\texttt{uint64}、\texttt{float32}、\texttt{float64}、\texttt{bool} & 标量，按值传 \\
\texttt{String} & 字符串（C 侧 \texttt{const char *}） \\
\texttt{T mut\&} / \texttt{T const\&} & 可变 / 只读引用（C 侧 \texttt{T *} / \texttt{const T *}） \\
\texttt{Option<\allowbreak{}T mut\&>\allowbreak{}} / \texttt{Option<\allowbreak{}T const\&>\allowbreak{}} & 可空引用（C 侧可以是空指针） \\
\texttt{T[\allowbreak{}]\allowbreak{}} + 另给长度 & 数组视图（C 侧 \texttt{T *}） \\
结构体 & 按值传，默认 C 布局 \\
\texttt{struct X packed \{ … \}} & 紧凑布局，对接 C 的 \texttt{\#pragma pack} \\
\bottomrule
\end{longtable}
}

\section{数组与「可空」：边界上的两个特殊话题}

数组用 \texttt{T[\allowbreak{}]\allowbreak{}} 传，长度单独给一个参数；如果返回的数组长度由某个形参决定，在末尾用 \texttt{[\allowbreak{}count]\allowbreak{}} 标注：

\par
\medskip
\begin{lstlisting}[style=chali]
#abi("C", "./native.dll", "packed_numbers")
Packed[] packed_numbers(int32 count) [count];
#abi("C", "./native.dll", "packed_fill")
void packed_fill(Packed[] values, int32 count);
\end{lstlisting}
\par
\medskip

「空」是这个边界上唯一允许出现空的地方。 ChaLi 内部的引用永远非空（第 7 章），但 C 函数可能返回空指针。这时候用 \texttt{Option<\allowbreak{}T mut\&>\allowbreak{}} / \texttt{Option<\allowbreak{}T const\&>\allowbreak{}} 接：

\par
\medskip
\begin{lstlisting}[style=chali]
#abi("C", "./native.dll", "packed_snapshot")
Option<Packed const&> packed_snapshot();
\end{lstlisting}
\par
\medskip

返回 \texttt{None} 就代表 C 那边给了个空指针；传参时也可以传\texttt{Option<\allowbreak{}Packet mut\&>\allowbreak{}:\allowbreak{}:\allowbreak{}None}，C 侧收到的就是空指针。

\begin{tip}{提示}
这是「空」唯一能跨过边界的通道。 跨过边界之后，一切又回到「引用非空」的老规矩。
\end{tip}

\texttt{packed} 结构体用于对接 C 侧 \texttt{\#pragma pack} 的紧凑布局——C 里用\texttt{\#pragma pack(\allowbreak{}push,\allowbreak{} 1)\allowbreak{}} 压过的地方，ChaLi 侧就声明成 \texttt{struct X packed \{ … \}}。

完整例子见 \href{https://github.com/rabbitc999/ChaLi/tree/master/doc/examples/abi_buffers}{\texttt{doc/\allowbreak{}examples/\allowbreak{}abi\_buffers/\allowbreak{}}}——它覆盖了结构体嵌套、紧凑布局、数组读写、可空引用、返回值长度等全部边界情况。

\section{宿主胶水：把 C++ 函数注册进来}

不想走动态库、想让 C++ 函数像普通函数一样调？写一个 \texttt{bind\allowbreak{}Cha\allowbreak{}Li\allowbreak{}Glue} 把它注册进你的宿主程序：

\par
\medskip
\begin{lstlisting}[style=cpp]
#include <cstdint>
#include <string>

struct ExportCounterTag;
using ExportCounter = chali_glue::Handle<ExportCounterTag, std::uint64_t>;

bool bindChaLiGlue(chali_glue::Registry &registry, Memory::String *error)
{
    return registry.bindTyped<std::uint32_t, Memory::String>(
               "demo_cpp_text_size",
               [](Memory::String arg0) {
                   return static_cast<std::uint32_t>(arg0.size());
               }, error);
}
\end{lstlisting}
\par
\medskip

ChaLi 侧就像声明普通函数：

\par
\medskip
\begin{lstlisting}[style=chali]
int32 demo_cpp_text_size(String text);

void main()
{
    println("{}", demo_cpp_text_size("hello"));
}
\end{lstlisting}
\par
\medskip

两条约定：

\begin{itemize}

\item 复杂宿主对象一律用整数句柄过边界——不要把 C++ 对象本体塞进 ChaLi 值里，而是用 \texttt{chali\_glue:\allowbreak{}:\allowbreak{}Handle<\allowbreak{}Tag,\allowbreak{} Abi\allowbreak{}Type>\allowbreak{}} 给它一个「名义类型」，ChaLi 侧拿到的是个整数；

\item 泛型函数实例化后键名带后缀：\texttt{demo\_cpp\_identity<\allowbreak{}int32>\allowbreak{}} 对应键名\texttt{demo\_cpp\_identity\_\_int32}。

\end{itemize}

宿主类型的名字是编译期就绑死的，所以这条路比 \texttt{abi} 强类型、报错也早，胶水编入 VM 宿主，并在运行程序前注册。

\subsection*{\texttt{impl}：把宿主函数绑定到协议}

结构体用继承列表声明普通接口实现；\texttt{impl 类型 :\allowbreak{} 接口 =\allowbreak{} 宿主绑定名;} 则把宿主函数接到单方法接口或转换协议。它不是一套任意 \texttt{impl Type \{ .\allowbreak{}.\allowbreak{}.\allowbreak{} \}} 语法。

下面以一组宿主绑定为例。这些绑定全部由你的宿主程序提供：使用胶水前必须先在 C++ 里注册同名函数，独立 AOT 产物不内置任何演示服务——只写声明并不能让它们存在，示例要在注册了这些绑定的宿主里运行才有意义。

\par
\medskip
\begin{lstlisting}[style=chali]
type HostHandle<T> = uint64;
impl uint64 : ExplicitConvert<HostHandle<int32>> = demo_counter_from_uint;
impl HostHandle<int32> : ExplicitConvert<uint64> = demo_counter_to_uint;

interface HostValue
{
    int64 hostValue();
}
impl HostHandle<int32> : HostValue = demo_counter_to_uint;
HostHandle<int32> demo_cpp_add(HostHandle<int32> left, HostHandle<int32> right);

void main()
{
    HostHandle<int32> left = cast<HostHandle<int32>>((uint64)20);
    HostHandle<int32> right = cast<HostHandle<int32>>((uint64)22);
    HostHandle<int32> total = demo_cpp_add(left, right);
    println("{} {}", cast<uint64>(total), total.hostValue());
}
\end{lstlisting}
\par
\medskip

接口签名、宿主参数和返回值必须一致。同步回调通过 \texttt{Native\allowbreak{}Call\allowbreak{}Context:\allowbreak{}:\allowbreak{}invoke} 与 \texttt{chali\_glue:\allowbreak{}:\allowbreak{}Function} 使用同一套 VM 规则；不要把函数值中的运行时地址当作宿主机器函数指针。完整胶水约定见\href{https://github.com/rabbitc999/ChaLi/blob/master/doc/08-host-integration.md}{宿主集成}。

\section{反方向：宿主调用 ChaLi}

宿主通过生成的 C ABI 头文件调用 AOT 原生动态库；先初始化，再调用导出函数，最后清理模块。宿主负责原生库与所借用资源的生命周期，不再通过 Machine::run 执行字节码。

在宿主原生函数内，可以通过 \texttt{chali\_glue:\allowbreak{}:\allowbreak{}Function} 或 \texttt{Native\allowbreak{}Call\allowbreak{}Context:\allowbreak{}:\allowbreak{}invoke} 同步调用 ChaLi 函数和闭包。它们支持高阶函数、引用和同步重入，并保留 VM 的错误传播与所有权约定。函数值中的运行时地址不能转换为宿主机器函数指针。

ChaLi 通过 VM 宿主加载胶水。完整接口和代码见\href{https://github.com/rabbitc999/ChaLi/blob/master/doc/08-host-integration.md}{宿主集成}及\href{https://github.com/rabbitc999/ChaLi/blob/master/example/native-callbacks/README.md}{VM 回调示例}。

\section{小结}

\begin{center}
\begin{tabular}{ll}
\toprule
你会了 & 关键点 \\
\midrule
三种集成 & \texttt{\#abi} 指令（轻）／宿主胶水（要编 C++）／VM 嵌入（宿主运行 ChaLi） \\
\texttt{\#abi} 指令 & \texttt{\#abi(\allowbreak{}"C",\allowbreak{} ".\allowbreak{}/\allowbreak{}x.\allowbreak{}dll",\allowbreak{} "C符号名")\allowbreak{} 返回类型 名字(\allowbreak{}形参)\allowbreak{};} \\
名字映射 & 第三个指令参数是 C 侧真实符号名，声明使用 ChaLi 本地名称 \\
结构体 & 按值传；ChaLi 侧声明字段一致的 \texttt{struct} \\
引用 & C 的 \texttt{T *} \ensuremath{\leftrightarrow} ChaLi 的 \texttt{T mut\&}，调用处写 \texttt{\&x} \\
字符串 & C 的 \texttt{const char *} \ensuremath{\leftrightarrow} ChaLi 的 \texttt{String} \\
数组 & \texttt{T[\allowbreak{}]\allowbreak{}} + 长度另给；返回长度用 \texttt{[\allowbreak{}count]\allowbreak{}} 标注 \\
「空」 & 只有 \texttt{Option<\allowbreak{}T \&>\allowbreak{}} 能把空送过边界 \\
紧凑布局 & \texttt{struct X packed \{ … \}} 对接 \texttt{\#pragma pack} \\
胶水 & \texttt{bind\allowbreak{}Cha\allowbreak{}Li\allowbreak{}Glue} 注册 C++ 函数；复杂对象用整数句柄 \\
反向 & VM 宿主执行，或胶水同步回调 fn \\
\bottomrule
\end{tabular}
\end{center}

\section{下一步}

你能调 C 了、能被 C 调了。最后一个技术问题是：怎么把程序交出去——变成对方机器上能跑的东西。

第 22 章讲部署原生产物：AOT 运行、检查配置、调试与交付。

翻到\hyperref[ch:22]{第 22 章 部署原生产物}。

\section{边界与验证}

ABI 一致只是互操作的前提之一。指针、字符串、回调上下文的拥有与有效期必须明确；同步回调走统一 VM 调用协议，运行时地址不是可直接执行的机器函数指针。

\section{习题}

前两题用于理解和基本实现，第三题检查边界，最后一题为综合挑战。先独立完成，再核对参考解答。

\begin{enumerate}

\item 为什么两边函数名相同仍可能无法安全调用？

\item 列出传入 UTF-8 字符串时宿主必须明确的三个约定。

\item 说明宿主保存一次同步调用的临时上下文为何可能出错。

\item 挑战：为一个宿主函数设计参数数目、返回类型和错误传播测试。

\end{enumerate}

\subsection*{参考解答与验收要点}

\begin{enumerate}

\item 还需要调用约定、参数宽度、布局、返回协议及生命周期一致。

\item 编码、长度或终止方式、缓冲区有效期与释放责任至少说明清楚。

\item 调用结束后上下文不再有效，后续使用不受原同步调用保证，需使用实际支持的长期注册协议。

\item 覆盖合法调用与不匹配调用，保留原始错误来源，检查没有把 VM 值地址强转成机器入口。

\end{enumerate}

\section{番外：借来的梨树与指针}

改写自《聊斋志异·种梨》。

有乡人货梨于市，颇甘芳，价腾贵。有道士求一梨，乡人叱之。旁人买而赠焉。道士食尽，留核，掘地埋之，俄而一树，枝叶扶疏，累累有梨，请众人任取。

乡人看得心热，忙将树的位置记在纸上，欲待散场，日日来摘。道士斫树而去，乡人依旧址伸手，只抓得一把尘土。乃怒曰：“这地址未变，梨何故不在？”

道士回身曰：“君记住地址，未曾取得所有权。对象既去，地址尚存，有何稀奇？”乡人仍不服，把旧指针换一个更大的整数，以为范围变宽，梨便能多活几天。

旁一木匠接过枝条，约定只在眼前加工。乡人又将这口头保证写成永久回调，待人散尽，遥唤木匠，哪还有上下文应答？

道士指着三个字教他：“谁创建，谁释放；借多久，何时用。跨语言的墙，挡不住生命周期的账。”

乡人回看梨车，梨已尽失，连车把也不知去了哪里。他捏着那张地址单，终于明白：ABI 对上，不过说明门牌格式认得；门后有没有人、谁有权搬家，仍须另外说清。

\subsection*{技术对照}

ABI 一致只是互操作的前提之一。指针、字符串、回调上下文的拥有与有效期必须明确；同步回调走统一 VM 调用协议，运行时地址不是可直接执行的机器函数指针。
